\documentclass[12pt,a4paper]{article}

\usepackage[T1]{fontenc}
\usepackage[utf8]{inputenc}
\usepackage[brazil]{babel}
\usepackage{graphicx}
\usepackage{geometry}
\usepackage{listings}
\usepackage{listingsutf8}
\usepackage{xcolor}
\usepackage{setspace}
\usepackage{float}
\usepackage{tikz}
\usepackage{tabularx}

\usetikzlibrary{positioning, shapes.geometric}
\onehalfspacing
\geometry{margin=2.5cm}

\lstset{
    language=SQL,
    basicstyle=\ttfamily\small,
    keywordstyle=\color{blue},
    commentstyle=\color{gray},
    stringstyle=\color{red},
    breaklines=true,
    showstringspaces=false
}

\newcommand{\montarinicio}[3]{
    % =========================
    % CAPA
    % =========================
    \begin{titlepage}
        \centering

        {\large Faculdade Anhanguera}\\[0.5cm]
        {\large Tecnólogo em DevOps}\\[3cm]

        {\Large \textbf{#1}}\\[2cm]

        {\large Gabriel Blank Stift Mousquer}\\[1cm]

        {\large Ivoti/RS}\\
        {\large 2026}

    \end{titlepage}

    % =========================
    % FOLHA DE ROSTO
    % =========================
    \begin{titlepage}
        \centering

        {\large Faculdade Anhanguera}\\[0.5cm]
        {\large Tecnólogo em DevOps}\\[2cm]

        {\Large \textbf{#1}}\\[1.5cm]

        \begin{flushright}
            Trabalho apresentado à disciplina de \\
            \textbf{#2} \\
            Professor: #3 \\
        \end{flushright}

        \vfill

        {\large Gabriel Blank Stift Mousquer}\\
        Matrícula: 2024201126\\

        \vfill

        {\large Ivoti/RS}\\
        {\large 2026}

    \end{titlepage}

    % =========================
    % SEÇÃO: TOC
    % =========================
    \tableofcontents
    \newpage
}


\begin{document}

\montarinicio
    {ESTRUTURAS DE CONTROLE E REPETIÇÃO}
    {Algoritmos e Programação Estruturada}
    {Anderson Emídio De Macedo Gonçalves}

\section{Codificação completa}

O código-fonte desenvolvido em linguagem C é apresentado a seguir.

\lstset{language=C}

\lstinputlisting[
    caption={Código em C},
    inputencoding=utf8/latin1,
    basicstyle=\ttfamily\scriptsize
]{./static/main.c}

\newpage

\section{Explicação do algoritmo}

O programa tem como objetivo ler números inteiros informados pelo usuário e calcular a soma desses valores até que seja inserido o número zero, utilizado como valor sentinela para encerrar o loop.

Inicialmente, são declaradas duas variáveis inteiras: \texttt{soma}, inicializada com zero, responsável por armazenar o total acumulado, e \texttt{valor}, inicializada com 1 para garantir que a condição do laço seja verdadeira na primeira verificação.

A estrutura de repetição utilizada é o laço \texttt{while}, que realiza o teste da condição antes de executar o bloco de código. A condição definida é \texttt{valor != 0}, ou seja, o loop continuará executando enquanto o valor informado pelo usuário for diferente de zero.

Dentro do laço, o programa solicita ao usuário a entrada de um número inteiro, que é armazenado na variável \texttt{valor}. Em seguida, esse valor é somado à variável \texttt{soma}, acumulando o resultado a cada iteração.

Quando o usuário insere o valor zero, a condição do \texttt{while} se torna falsa na próxima verificação, encerrando o loop. O valor zero também é somado, porém não altera o resultado final.

Após o término do laço, o programa exibe o valor total da soma acumulada.

A implementação garante que:
\begin{itemize}
    \item O laço execute com teste de condição no início, conforme proposto;
    \item O usuário possa inserir múltiplos valores;
    \item O loop seja encerrado corretamente com o valor sentinela;
    \item A soma final seja exibida corretamente.
\end{itemize}

\newpage

\section{Execução do código}

Foram realizados testes com diferentes conjuntos de valores para validar o funcionamento do programa.

\subsection{Caso 1: 4, 3, 2, 1, 0}

\begin{figure}[H]
    \centering
    \includegraphics[width=0.8\textwidth]{./static/img-1.png}
\end{figure}

\subsection{Caso 2: 0}

\begin{figure}[H]
    \centering
    \includegraphics[width=0.8\textwidth]{./static/img-2.png}
\end{figure}

\subsection{Caso 3: 5, 4, -15, 0}

\begin{figure}[H]
    \centering
    \includegraphics[width=0.8\textwidth]{./static/img-3.png}
\end{figure}

\end{document}
